%%
%% Ergänzung zu Kapitel 1: Grundlagen und verwandte Arbeiten
%%

\begin{neu}
\section{Studienintervention und KI-Prototyp}

Die Arbeit verbindet drei technisch zusammenhängende, wissenschaftlich jedoch getrennt zu behandelnde Ebenen: die Erzeugung visueller Reizmaterialien, die kontrollierte Durchführung der App-Studie und die Evaluation einer später möglichen automatischen Reizerkennung. Die produktive Studie operationalisiert den von Tabibnia et al. untersuchten Vergleich zwischen Cue-Matching und Cue-Labeling. Dessen Interpretation setzt voraus, dass die Bedingungen bis auf die Art der Zuordnungsaufgabe möglichst vergleichbar bleiben. Arbeiten zur virtuellen und augmentierten Cue-Exposition zeigen außerdem, dass die Gestaltung der Reize und der Nutzungskontext sowohl die ausgelöste Reaktion als auch die Übertragbarkeit gelernter Effekte beeinflussen können \cite{tabibnia_cuelabeling_2026,pericot_valverde_vr_smoking_cessation_2019,dougan_ar_smoking_cessation_2025}. CueLens verwendet in der Wirksamkeitsstudie deshalb statisch gebündelte Bilder und festgelegte Antwortoptionen, sodass alle Teilnehmenden einen reproduzierbaren Aufgabenbestand erhalten.

Für die Erstellung eines Teils der Cue-Bilder wurden generative Diffusionsmodelle eingesetzt. Denoising Diffusion Probabilistic Models lernen einen schrittweisen Prozess, der aus verrauschten Repräsentationen wieder Bilddaten erzeugt \cite{ho_diffusion_2020}. Latent-Diffusion-Modelle verlagern diesen rechenintensiven Prozess in den komprimierten Repräsentationsraum eines Autoencoders und ermöglichen über Konditionierung, beispielsweise durch Texteingaben, die Erzeugung hochaufgelöster Bilder. Die technische Grundlage von Stable Diffusion gehört zu dieser Modellklasse \cite{rombach_latent_diffusion_2022}. In CueLens dienen die generativen Modelle ausschließlich der vorgelagerten Erstellung möglicher Bildressourcen. Die erzeugten Bilder werden manuell gesichtet, ausgewählt und anschließend unverändert in die App übernommen; während einer Studiensituation wird kein Bild neu generiert. Damit wird die stochastische Bildsynthese in einen festen und für alle Teilnehmenden gleichen Reizbestand überführt.

Davon getrennt untersucht der KI-Prototyp die Klassifikation bereits vorhandener Bilder. Bei einem überwachten Klassifikationsverfahren wird aus gelabelten Beispielen eine Zuordnung von Eingaben zu Klassen gelernt. Entscheidend ist dabei nicht allein der Fehler auf den Trainingsdaten, sondern die Vorhersagegüte bei zuvor ungesehenen Bildern \cite{james_islr_2021}. Transfer Learning nutzt hierfür Repräsentationen, die zunächst auf einer umfangreicheren Ausgangsaufgabe gelernt wurden, und passt sie an eine kleinere Zielaufgabe an. Dadurch kann der Bedarf an projektspezifisch gelabelten Daten sinken; der Erfolg hängt jedoch davon ab, wie gut Ausgangs- und Zieldomäne zueinander passen \cite{pan_transfer_learning_2010}. Für den Prototyp wurden mit MobileNetV3Small und EfficientNet-Lite0 zwei Modellfamilien ausgewählt, deren Entwurf ausdrücklich auf ein günstiges Verhältnis von Vorhersagegüte und Ressourcenbedarf ausgerichtet ist. MobileNetV3 wurde mittels hardwarebezogener Architektursuche für mobile Prozessoren optimiert, während EfficientNet Tiefe, Breite und Eingabeauflösung eines Netzes gemeinsam skaliert \cite{howard_mobilenetv3_2019,tan_efficientnet_2019}.

Die Trennung von Studienintervention und Klassifikationsprototyp ist eine Voraussetzung für die interne Nachvollziehbarkeit der Wirksamkeitsstudie. Würde ein Klassifikationsmodell während der Studiensituationen Reize, Vergleichsbilder oder Wörter auswählen, würden seine Fehler zu einem Bestandteil der experimentellen Manipulation. Ein falsch negatives Ergebnis könnte einen tatsächlich vorhandenen Rauchreiz übersehen. Ein falsch positives Ergebnis könnte für ein neutrales Bild eine rauchbezogene Antwortoption erzeugen. Präzision und Recall beschreiben unterschiedliche Seiten dieser Fehlentscheidungen und müssen deshalb getrennt bewertet werden \cite{manning_ir_2008,james_islr_2021}. Ein Unterschied der Craving-Werte ließe sich bei einer solchen Kopplung nicht mehr eindeutig auf Cue-Labeling zurückführen, sondern könnte zusätzlich von Modellgüte, Klassenschwellen und Zusammensetzung der verwendeten Bilder abhängen.

Auch die klinische Einordnung spricht gegen eine vorschnelle Gleichsetzung technischer Funktionsfähigkeit mit einer wirksamen Intervention. Die WHO nennt Smartphone-Apps und KI-basierte Software als mögliche digitale Unterstützung der Tabakentwöhnung, bewertet die Evidenz hierfür jedoch nur mit geringer Gewissheit und fordert eine sorgfältige Gestaltung, fortlaufende Überwachung und Evaluation digitaler Angebote \cite{who_tobacco_cessation_guideline_2024}. Der KI-Prototyp liefert daher zunächst einen technischen Machbarkeits- und Modellvergleich. Er erbringt keinen eigenständigen Nachweis, dass eine automatisierte Reizerkennung Rauchverlangen reduziert oder den Rauchstopp unterstützt.

Die lokale Ausführung der Klassifikation hat gegenüber einem serverseitigen Bilddienst einen datenschutzbezogenen Vorteil: Bilder können auf dem Android-Gerät decodiert und ausgewertet werden, ohne dass die Dateien oder daraus erzeugte Merkmalsvektoren an einen externen Analysedienst übertragen werden. Die BSI-Richtlinie für mobile Gesundheitsanwendungen fordert, Architektur, Datenspeicherung, Netzwerkkommunikation und den Schutz sensibler Datenelemente gemeinsam zu betrachten \cite{bsi_tr03161_mobile_2024}. On-Device-Inferenz reduziert damit einen Datenfluss, beseitigt aber weder methodische Anforderungen an die Modellvalidierung noch urheberrechtliche und persönlichkeitsrechtliche Fragen des Bildmaterials.

Für die vorliegende Arbeit werden die Resultate der App-Studie und der KI-Evaluation deshalb getrennt berichtet. Die App-Studie beantwortet die verhaltensmedizinische Frage, ob sich der im Labor beobachtete Unterschied zwischen Cue-Matching und Cue-Labeling mit kontrollierten mobilen Aufgaben im Alltag nachvollziehen lässt. Die KI-Evaluation prüft dagegen, ob ein lokales Klassifikationsmodell technisch als Grundlage für eine spätere kontextbezogene Reizauswahl geeignet sein könnte. In der aktuellen Wirksamkeitsstudie werden keine selbst aufgenommenen Bilder klassifiziert. Eine spätere Verbindung beider Komponenten wäre ein neuer Untersuchungsgegenstand und müsste hinsichtlich Modellfehlern, Vergleichbarkeit der Reize, Datenschutz und möglicher zusätzlicher Belastungen erneut geplant und evaluiert werden.
\end{neu}
